\section*{Chapitre \ref{ch:iv2:crypto-for-traders-and-researchers} --- Crypto pour traders et chercheurs}
\markboth{\protect\omTocMark{}{Réponses~: Crypto}}{\protect\omTocMark{}{Réponses~: Crypto}}

\begin{solution}{iq:iv2:crypto-for-traders-and-researchers:1}
On paie 100~\$ par intervalle, 300~\$ par jour et environ 109\,500~\$ par an, soit 10.95~\%. Quand le perpétuel se traite au-dessus du spot, les longs paient en général les shorts.

$0.0001 \times \$1$ million $= \$100$, trois fois par jour, 365 jours par an.

Le funding sert à maintenir un contrat sans échéance proche du spot. Si le perpétuel est cher, le détenir à l'achat coûte de l'argent, donc les longs sortent et les shorts entrent.

\iqwrong{Traiter 0.01~\% comme négligeable parce que le chiffre paraît petit. Sur 1\,095 intervalles par an, il se compose en un vrai carry.}
\end{solution}

\begin{solution}{iq:iv2:crypto-for-traders-and-researchers:2}
Cela fait environ 10.1~\% annualisé. On achète le spot et on vend le contrat à terme, ce qui verrouille 1\,500~\$ par bitcoin, soit 2.5~\% sur les 90 jours.

$1\,500/60\,000 = 2.5~\%$, et $2.5~\% \times 365/90 = 10.14~\%$. À l'échéance, le contrat à terme converge vers le spot, donc les deux jambes verrouillent la différence.

\iqwrong{Comparer les 2.5~\% à un taux annuel sans les annualiser. De plus, il faut du capital sur les deux jambes~: le rendement réel sur le capital est inférieur à la base.}
\end{solution}

\begin{solution}{iq:iv2:crypto-for-traders-and-researchers:3}
On gagne 3.33 bitcoins à 60\,000 et on perd 5 bitcoins à 40\,000. En dollars, cela fait 200\,000~\$ gagnés et 200\,000~\$ perdus.

$10^6(1/50\,000 - 1/60\,000) = 3.33$, et $10^6(1/50\,000 - 1/40\,000) = -5$.

Le P\&L en dollars est le P\&L en coins multiplié par le prix courant~: $N(P_1/P_0 - 1)$, qui est linéaire. Le P\&L en coins, lui, ne l'est pas~: on perd davantage de coins à la baisse qu'on n'en gagne à la hausse.

\iqwrong{Supposer qu'il se comporte comme un contrat qui paie en dollars. Le P\&L en dollars concorde bien, mais on le détient en coins dont la valeur a changé, et la marge évolue elle aussi.}
\end{solution}

\begin{solution}{iq:iv2:crypto-for-traders-and-researchers:4}
Vers 45\,250.

La marge vaut 10~\% du notionnel, soit 5\,000 par bitcoin. La liquidation intervient lorsque les fonds propres tombent au niveau de maintenance de 0.5~\%, soit 250 par bitcoin. Le prix peut donc baisser de 4\,750, jusqu'à 45\,250.

\iqwrong{Répondre 45\,000. C'est le prix de faillite, où les fonds propres sont nuls. La liquidation doit intervenir avant.}
\end{solution}

\begin{solution}{iq:iv2:crypto-for-traders-and-researchers:5}
Cela coûte environ 22\,222 USDC, soit un prix moyen de 2\,222 par ETH. Le prix marginal après la transaction est d'environ 2\,469.

L'invariant est $k = 100 \times 200\,000 = 2 \times 10^7$. Après la transaction, $x = 90$, donc $y = 222\,222$. Le prix marginal est $y/x = 2\,469$.

\iqwrong{Payer le spot de 2\,000 sur les 10 unités. Le prix se déplace le long de la courbe à mesure qu'on achète, et il faut le payer~: 11~\% en moyenne pour 10~\% du pool.}
\end{solution}

\begin{solution}{iq:iv2:crypto-for-traders-and-researchers:6}
On achète à 0.995 et on rachète à 0.999, ce qui rapporte 0.4~\% par manche. Les délais, les limites de rachat ou les doutes sur la réserve peuvent empêcher l'écart de se refermer. À 1.01, on émettrait à 1.00 et on vendrait.

Le rachat donne $1 - 0.001 = 0.999$. Chaque aller-retour réduit aussi l'offre, ce qui aide à ramener le prix vers la parité.

\iqwrong{Supposer que n'importe qui peut le faire. En général, seuls les comptes agréés peuvent émettre et racheter directement, de sorte que l'arbitrage dépend d'eux, et un ancrage peut tenir ou céder selon leur appétit.}
\end{solution}

\begin{solution}{iq:iv2:crypto-for-traders-and-researchers:7}
Non~: c'est à peu près le seuil de rentabilité. Le gain brut est de 60~\$, soit 10 points de base, et les frais absorbent environ 59.97~\$.

$60\,000 \times 0.0005 + 59\,940 \times 0.0005 = 59.97$. Le gain net est donc de 3 centimes sur un bitcoin.

\iqwrong{Comparer les prix et s'arrêter là. Un écart visible doit couvrir les deux jeux de frais, et il peut avoir disparu avant que les deux jambes soient exécutées.}
\end{solution}

\begin{solution}{iq:iv2:crypto-for-traders-and-researchers:8}
Votre position vaut environ 5.72~\% de moins que la détention. C'est le même écart de 5.72~\% si le prix est divisé par deux.

La valeur du pool est $V = 2\sqrt{kP}$, et la détention vaut $(P/P_0 + 1)V_0/2$. Pour un rapport de prix $r$, le ratio vaut donc $2\sqrt{r}/(1 + r)$. En $r = 2$, il vaut $2\sqrt{2}/3 = 0.9428$. En $r = 1/2$, c'est la même valeur.

\iqwrong{Croire que le pool perd de l'argent en dollars quand le prix double. Le pool a gagné environ 41~\%. Il est seulement inférieur à ce que procure la détention, et cela se produit avant les frais.}
\end{solution}

\begin{solution}{iq:iv2:crypto-for-traders-and-researchers:9}
On gagne environ 9.1~\% sur le capital avec une marge de 20~\%, et 7.3~\% à 50~\%. Un funding qui devient négatif ou un rebond brutal qui vous force à sortir met fin à l'opération.

Le funding vaut $0.0001 \times 1\,095 = 10.95~\%$ sur le notionnel. Le capital est constitué du million de dollars de spot et de 200\,000~\$ de marge, donc $10.95/1.2 = 9.1~\%$. Avec 500\,000~\$ de marge, $10.95/1.5 = 7.3~\%$.

Il y a un arbitrage à faire~: une marge plus faible rapporte un peu plus, mais un rebond de 20~\% effacerait alors la marge du short, alors même que la jambe spot a gagné le même montant sur une autre plateforme.

\iqwrong{Donner 10.95~\% comme rendement. C'est le rendement sur le notionnel. Et le gain sur le spot ne peut pas servir de marge, sauf si les deux jambes partagent une même plateforme avec une marge croisée.}
\end{solution}

\begin{solution}{iq:iv2:crypto-for-traders-and-researchers:10}
On est long le contrat à terme et short le perpétuel. Sur 90 jours, on encaisse 2.7~\% de funding et on abandonne 2.5~\% de base, soit 0.2~\%, ou 2\,000~\$ sur un million, si le funding se maintient.

Le funding cumulé vaut $0.0001 \times 270 = 2.7~\%$.

Le spot n'intervient pas, car à l'échéance le contrat à terme vaut le spot et le perpétuel est proche du spot. Le gain de la position longue sur le contrat à terme est $S - F_0$, et celui de la position courte sur le perpétuel est $P_0 - S$. Ensemble, cela fait $P_0 - F_0$, soit moins la base, sans $S$.

Le risque tient à l'évolution du funding. On reçoit un taux variable et on a fixé une base. C'est donc un swap~: un funding variable contre une base fixe. La marge sur les deux jambes et l'écart du perpétuel par rapport au spot à l'échéance complètent la liste.

\iqwrong{Dire qu'il s'agit d'un arbitrage. Ce serait le cas si le funding était fixe. Or il est descendu bien en dessous de ce niveau sur des marchés calmes.}
\end{solution}

\begin{solution}{iq:iv2:crypto-for-traders-and-researchers:11}
Votre position vaut 60\,000~\$ dans les deux cas. C'est une vraie position en dollars. Le P\&L en coins de la position longue est concave~: il gagne moins à la hausse qu'il ne perd à la baisse.

À 30\,000, le short inverse rapporte $-60\,000(1/60\,000 - 1/30\,000) = +1$ bitcoin. On a donc 2 bitcoins à 30\,000, soit 60\,000~\$. À 90\,000, cela donne $-60\,000(1/60\,000 - 1/90\,000) = -1/3$, soit $2/3$ de bitcoin à 90\,000, soit 60\,000~\$.

Pour une position longue avec une marge en coins, une baisse vous coûte à la fois davantage de coins et réduit la valeur de ceux qui restent. La marge est donc consommée deux fois plus vite, ce qui avance la liquidation par rapport à un contrat margé en dollars.

\iqwrong{Croire qu'une position longue margée en coins et une position longue margée en dollars ont le même risque de liquidation. Pour le même notionnel en dollars et la même marge, la version en coins est plus fragile à la baisse.}
\end{solution}

\begin{solution}{iq:iv2:crypto-for-traders-and-researchers:12}
Le prix finit en baisse de 12.5~\%. Si les liquidations représentent 50 millions de dollars par point de pourcentage, la chute n'a pas de limite dans ce modèle.

Chaque point de baisse force 30 millions de dollars de ventes, et cela fait baisser le prix de $30/50 = 0.6~\%$ supplémentaires. Le total vaut donc $5~\% \times (1 + 0.6 + 0.36 + \dots) = 5~\%/(1 - 0.6) = 12.5~\%$.

Pour un rapport égal à 1, la série ne converge pas. Le prix baisse donc jusqu'à épuisement des positions ou du levier, ou jusqu'à la suspension de la cotation.

\iqwrong{Croire que 5~\% est l'événement. Ce qui compte est le multiplicateur de cascade, $1/(1 - b)$, et $b$ reste caché tant que le levier ne s'est pas accumulé.}
\iqfollow{Que mesureriez-vous à l'avance~?}{Le nombre de positions ouvertes avec leurs niveaux de liquidation estimés, et la profondeur du carnet. Un rapport élevé du premier sur la seconde est un signal d'alerte.}
\end{solution}

\begin{solution}{iq:iv2:crypto-for-traders-and-researchers:13}
Le fonds d'assurance absorbe le déficit de 10\,000 dollars. Pour une exécution à 45\,250, les 5\,000 dollars de marge restante reviennent en général au fonds. Si le fonds est vide, la plateforme réduit les positions opposées rentables, c'est ce qu'on appelle l'auto-deleveraging (ADL).

Sous le prix de faillite de 45\,000, la position a une valeur nette négative. À 44\,500, cela fait $500 \times 20 = \$10\,000$. À 45\,250, la valeur nette restante vaut $250 \times 20 = \$5\,000$.

Les traders gagnants peuvent donc voir leurs couvertures réduites au pire moment. Un teneur de marché qui se croyait neutre peut se retrouver exposé directionnellement. C'est un risque réel qu'il faut intégrer dans le prix.

\iqwrong{Supposer qu'un short couvert ne peut pas être touché par les liquidations des autres. L'auto-deleveraging peut fermer votre jambe rentable et laisser votre autre jambe exposée.}
\end{solution}

\begin{solution}{iq:iv2:crypto-for-traders-and-researchers:14}
Cela coûte environ 22\,289 USDC, avec des frais d'environ 67. La bande va d'environ 1\,994 à 2\,006.

Les frais étant prélevés sur l'entrée, $0.997\,\mathrm{d}y$ entre dans la courbe. Donc $\mathrm{d}y = 22\,222/0.997 = 22\,289$, et les frais valent 66.87.

Un acheteur paie $2\,000/0.997 = 2\,006$ à la marge, et un vendeur reçoit $2\,000 \times 0.997 = 1\,994$. Dans cette bande, personne n'a donc intérêt à trader contre le pool.

Le pool joue donc le rôle d'un teneur de marché qui affiche une demi-fourchette fixe égale aux frais.

\iqwrong{Penser que les frais ne comptent qu'à la fin. Ils décalent le point où l'arbitrage commence, si bien que le prix d'un pool peut s'écarter de l'équilibre d'un montant égal aux frais, sans que les arbitragistes le corrigent.}
\end{solution}

\begin{solution}{iq:iv2:crypto-for-traders-and-researchers:15}
Acheter de l'ETH avec des dollars, le revendre contre du BTC à 0.0340, puis revendre le BTC contre des dollars. Cela rapporte environ 2.0~\% brut, et environ 1.69~\% après frais. Il faudrait environ 66 points de base par jambe pour l'annuler.

Le taux implicite vaut $2\,000/60\,000 = 0.03333$, et $0.0340/0.03333 = 1.02$. Donc 1~\$ devient 1.02~\$ avant frais, et $1.02 \times 0.999^3 = 1.0169$ après.

Les frais d'annulation résolvent $(1 - f)^3 = 1/1.02$~: $f = 0.66~\%$.

\iqwrong{Faire la boucle dans le mauvais sens~: $1/1.02$ perd 2~\%. Il faut repérer la jambe chère. Ici, l'ETH est cher par rapport au BTC, donc on vend l'ETH contre du BTC.}
\end{solution}

\begin{solution}{iq:iv2:crypto-for-traders-and-researchers:16}
Il vaut environ 0.96. Un prix inférieur à 0.92 signifie que le marché doute des 92~\% restants, ou qu'il s'agit d'une décote de liquidité. Il faut dimensionner la position pour le cas où l'on ne pourrait pas sortir pendant des semaines.

$0.92 + 0.08 \times 0.5 = 0.96$. Même sans aucune récupération des 8~\% bloqués, la pièce vaut 0.92. Un prix de 0.88 ne peut donc pas venir de ces seuls 8~\%.

Ce prix traduit l'un de trois cas~: les autres réserves sont elles aussi en doute, les rachats sont suspendus et les détenteurs ont besoin de dollars tout de suite, ou il existe une prime de contagion. Le gain potentiel est donc d'environ 9~\% si vous avez raison.

\textbf{Dimensionnement.} Il faut supposer qu'aucun rachat n'ait lieu pendant un certain temps, et calculer la perte à un prix bien plus bas. Il faut aussi se demander si l'on peut tenir la position ou si le capital est nécessaire ailleurs.

\iqwrong{Prendre la seule réserve bloquée pour réponse. Une ruée est un problème de calendrier et d'accès au rachat, autant que de valeurs.}
\end{solution}

\begin{solution}{iq:iv2:crypto-for-traders-and-researchers:17}
Environ 200 millions de jetons sont émis, le prix passe de 1.00 à 0.25, et le prix moyen réalisé vaut 0.50. La conception transforme donc la vente en davantage de vente.

Vendre des jetons fait sortir des dollars du pool. Pour obtenir 100 millions de dollars, le côté dollar doit passer de 200 à 100. Le côté jeton passe donc de 200 à 400, ce qui exige 200 millions de jetons. Le prix vaut $100/400 = 0.25$.

\textbf{Le problème.} La garantie est le jeton frère, et ce sont les attaques par rachat qui en font chuter le prix. Chaque dollar racheté émet davantage de jetons à un prix plus bas. Une fois la ruée engagée, il faut donc plus d'offre pour chaque dollar, dans une boucle sans plancher. Rien d'extérieur au système ne la soutient.

\iqwrong{Compter 100 millions de jetons parce que le prix vaut 1. Cela ignore l'impact de sa propre vente, et cet impact double le nombre de jetons nécessaires.}
\end{solution}

\begin{solution}{iq:iv2:crypto-for-traders-and-researchers:18}
Le funding vaut 0.01~\%, 0.01~\%, 0.05~\%, 0.01~\% et $-0.15~\%$.

À l'intérieur de la plage de $P$, le clamp annule la prime, et $F = I$. Au-delà, $F$ suit $P$ avec un décalage de 0.05~\%~:
\begin{itemize}
\item $P = 0.02$~: $I - P = -0.01$, donc $F = 0.01$.
\item $P = 0.06$~: $I - P = -0.05$, donc $F = 0.01$.
\item $P = 0.10$~: $I - P = -0.09$, borné à $-0.05$, donc $F = 0.05$.
\item $P = -0.03$~: $I - P = 0.04$, donc $F = 0.01$.
\item $P = -0.20$~: $I - P = 0.21$, borné à $+0.05$, donc $F = -0.15$.
\end{itemize}

Le funding reste donc égal au terme d'intérêt pour un $P$ compris entre $-0.04~\%$ et $0.06~\%$. C'est pourquoi le funding affiche souvent une seule et même valeur. Il agit comme une zone morte avec une pente linéaire à l'extérieur.

\begin{dated}{2026-11}{Une règle publiée}\label{dat:iv2:crypto-for-traders-and-researchers:funding}
Une grande plateforme publie cette formule, avec un terme d'intérêt de 0.01~\% par intervalle sur la plupart des contrats, et trois intervalles par jour. Les formules et les intervalles diffèrent selon les plateformes et les contrats. Il faut consulter la page en vigueur avant de trader.
\end{dated}

\iqwrong{Prendre 0.01~\% pour un signal de faible demande sur le marché. Dans la zone morte, ce chiffre ne dit rien d'autre que la faiblesse de la prime.}
\end{solution}

\begin{solution}{iq:iv2:crypto-for-traders-and-researchers:19}
Le coût espéré vaut 0.9~\% du capital par an. Répartir sur trois plateformes laisse ce coût inchangé et ramène la perte due à une seule défaillance de 18~\% à 6~\%, mais la défaillance des trois coûterait toujours 18~\%. Un portefeuille couvert perd aussi sa couverture.

$0.3 \times 0.05 \times 0.6 = 0.9~\%$. Avec trois plateformes détenant chacune 10~\%, le coût espéré vaut $3 \times 0.1 \times 0.05 \times 0.6 = 0.9~\%$. Une défaillance isolée coûte désormais 6~\%.

Un portefeuille acheteur en spot sur une plateforme et vendeur en perpétuel sur une autre n'est neutre que tant que les deux existent. Quand l'une fait défaut, l'autre jambe devient soudainement directionnelle. Et la défaillance survient souvent dans un marché violent, si bien qu'on peut perdre beaucoup plus sur cette jambe que sur les fonds gelés.

\begin{dated}{2026-11}{Les chiffres d'une défaillance}\label{dat:iv2:crypto-for-traders-and-researchers:fail}
Les actes de faillite d'une grande plateforme offshore défaillante en novembre 2022 font apparaître environ 11.2 milliards de dollars dus aux clients de la plateforme internationale, pour environ 2.5 milliards de dollars d'actifs à la date de la requête.
\end{dated}

\iqwrong{Prendre la diversification pour une élimination du risque de plateforme. Elle réduit la taille de chaque perte, mais la perte espérée reste la même, sauf si les plateformes diffèrent en qualité.}
\end{solution}

\begin{solution}{iq:iv2:crypto-for-traders-and-researchers:20}
$F = S\kappa/(\kappa - r)$. Pour $\kappa = 1\,095$ et $r = 5~\%$, $F$ dépasse le spot d'environ 0.46 point de base, et le funding vaut environ 5~\% par an~: c'est le taux d'intérêt.

\textbf{Démonstration.} On cherche $F = cS$. On détient $c$ unités de spot, financées au taux $r$, et on vend 1 perpétuel. Sur $\mathrm{d}t$, le spot et le perpétuel se compensent. Les intérêts coûtent $rcS\,\mathrm{d}t$, et l'on reçoit le funding $\kappa(F - S)\,\mathrm{d}t = \kappa(c - 1)S\,\mathrm{d}t$. Sans risque, il n'y a pas de profit~: $\kappa(c - 1) = rc$. Donc $c = \kappa/(\kappa - r)$.

$c - 1 = 0.05/1\,094.95 = 4.57 \times 10^{-5}$. Le funding vaut $\kappa(c - 1) = rc \approx 5~\%$.

Le funding est donc le coût de portage, et la prime sert à payer ce coût. Avec un rendement $q$ sur la pièce, il vaut $\kappa - (r - q)$. Quand $\kappa \to \infty$, $F \to S$.

\iqwrong{Penser que le funding doit tomber à zéro sans prime. Une prime nulle avec un taux positif donnerait un portage gratuit, puisqu'un acheteur ne paie pas d'intérêt et supporte toute l'exposition.}
\iqfollow{Pourquoi le funding réel diffère-t-il de cette formule~?}{La demande de levier et le coût des marges le poussent à la hausse, et le capital entre plateformes ne circule pas librement. La formule est l'ancrage sans frictions, et elle donne le signe de l'excès.}
\end{solution}

\begin{solution}{iq:iv2:crypto-for-traders-and-researchers:21}
Environ 68~\% à 10x, 43~\% à 5x et 7.7~\% à 2x, selon la règle de la barrière.

Sans dérive, la probabilité qu'un log-prix touche une barrière $b$ d'ici à $T$ vaut $2\Phi(-b/(\sigma\sqrt{T}))$. Ici, $\sigma\sqrt{T} = 0.8\sqrt{30/365} = 0.229$. Donc $b = \ln 1.1$, $\ln 1.2$ et $\ln 1.5$ donnent $2\Phi(-0.415) = 0.678$, $2\Phi(-0.795) = 0.427$ et $2\Phi(-1.77) = 0.077$.

Un passage de levier de 5x à 2x réduit la probabilité d'environ cinq fois et demie. Passer de 10x à 5x ne la réduit que d'environ 1.6 fois. Le risque est donc très non linéaire en fonction de la distance à la barrière.

\iqwrong{Utiliser $\Phi$ de la distribution terminale, ce qui diviserait ces chiffres par deux. Ce qui compte, c'est le toucher, puisque la position est liquidée dès que la barrière est atteinte, même si le prix revient ensuite.}
\iqfollow{Que vous donne une simulation~?}{Un peu moins avec un contrôle quotidien, car le suivi discret manque certains passages. Les moteurs de liquidation réels surveillent en continu, il faut donc retenir le chiffre continu.}
\end{solution}

\begin{solution}{iq:iv2:crypto-for-traders-and-researchers:22}
On gagne environ 20 points de base de capital par jour, soit environ 73~\% par an. Chaque tranche de 1 million de dollars déplacée a un écart type d'environ 6\,000~\$. On l'élimine en couvrant le mouvement sur un perpetual, ce qui coûte des frais et un peu de base.

Le revenu vaut $\$5 \text{ million} \times 0.001 = \$5\,000$. Les transferts coûtent $\$2 \text{ million} \times 0.0005 = \$1\,000$. Donc $\$4\,000/\$2 \text{ million} = 0.2~\%$ par jour.

Le risque vaut $0.8\sqrt{0.5/8\,760} = 0.60~\%$ sur chaque tranche de 1 million de dollars tant que la pièce est en transit, soit 6\,000~\$. C'est un bénéfice d'une journée et demie sur un mouvement de 1 million de dollars s'il va contre soi à 1 écart type. On vend un perpetual à découvert pendant que la pièce est en transit, puis on procède au débouclage à l'arrivée. Cela ajoute deux frais supplémentaires, disons 5 points de base chacun~: 1\,000~\$ par tranche de 1 million de dollars déplacée, soit 2\,000~\$ par jour. Le gain net tombe à 2\,000~\$, soit 10 points de base de capital par jour, soit environ 37~\% par an.

Le 73~\% n'est pas extensible. Un avantage de ce type est vite éliminé par la concurrence. Doubler le capital ne double pas une opportunité de 5 millions de dollars, et à mesure que de plus en plus d'entreprises la poursuivent, les 10 points de base se réduisent.

\iqwrong{Extrapoler le rendement à n'importe quel capital. L'arbitrage a une capacité limitée. Et la défaillance d'une plateforme coûterait davantage que le profit d'une année.}
\end{solution}

\begin{solution}{iq:iv2:crypto-for-traders-and-researchers:23}
$V = 2\sqrt{kP}$ a un delta égal à $\sqrt{k/P}$, et $\Gamma = -\tfrac12\sqrt{k}\,P^{-3/2}$. La perte couverte vaut $\tfrac12|\Gamma|\sigma^2P^2\,\mathrm{d}t = V\sigma^2/8\,\mathrm{d}t$. À 80~\%, cela fait 8~\% par an, et avec la décroissance de la valeur, environ 7.7~\% en espérance sur un an.

\textbf{Démonstration.} $\tfrac12 \times \tfrac12\sqrt{k}P^{-3/2} \times \sigma^2P^2 = \tfrac14\sqrt{kP}\,\sigma^2 = (V/2)\sigma^2/4 = V\sigma^2/8$.

Pour une position couverte en delta, $\mathrm{d}V - \Delta\,\mathrm{d}P = \tfrac12\Gamma\sigma^2P^2\,\mathrm{d}t$. Il s'agit donc d'un pur coût de gamma. Comme $V$ a une dérive de $-\sigma^2/8$, l'espérance de la valeur après un an vaut $1 - e^{-0.08} = 7.69\%$. Une simulation confirme ce résultat. Cette perte s'appelle \index{loss-versus-rebalancing}loss-versus-rebalancing.

\textbf{Pourquoi elle ne disparaît pas.} La couverture est rééquilibrée en continu. Chaque mouvement du pool vend bas et achète haut face à l'arbitragiste, et ces pertes sont figées. La perte impermanente, mesurée par rapport à la simple détention, peut revenir à zéro. Ce coût, lui, ne revient jamais à zéro.

\iqwrong{Traiter les revenus de frais du LP comme un profit. Ces frais rémunèrent la position vendeuse de gamma, et il faut les comparer à ce coût courant.}
\end{solution}

\begin{solution}{iq:iv2:crypto-for-traders-and-researchers:24}
On perd 20\,000~\$ à 2\,420. C'est exactement la perte impermanente par rapport à la détention. Quand le prix revient, le P\&L retombe à zéro, sans aucune transaction de couverture~; il s'agissait donc d'une perte véritablement indépendante du chemin suivi.

$k = 2 \times 10^9$. À 2\,420, $V = 2\sqrt{kP} = 2 \times 2.2 \times 10^6 = 4.4$ millions de dollars. Elle a gagné 400\,000~\$. La position vendeuse perd $1\,000 \times 420 = \$420\,000$. Net~: $-\$20\,000$.

Détenir les actifs aurait valu $1\,000 \times 2\,420 + 2\,000\,000 = 4.42$ millions de dollars. La perte impermanente est donc la même, soit 20\,000~\$, ou $2\sqrt{1.21}/2.21 - 1 = -0.45~\%$.

Une couverture statique au départ donne donc exactement la perte impermanente~: elle est indépendante du chemin. Une couverture continue donne une perte réalisée qui dépend du chemin, de l'ordre de $V\sigma^2/8\,\mathrm{d}t$, et qui s'accumule. Pour un grand mouvement dans un seul sens, la couverture statique perd davantage. Pour un aller-retour volatil, la statique ne perd rien, tandis que la couverture continue a payé chaque oscillation.

C'est un choix sur le gamma~: la couverture statique est exposée au niveau final, alors que la couverture continue paie la variance réalisée.

\iqwrong{Croire que la perte impermanente et le loss-versus-rebalancing sont la même chose. Elles sont égales en moyenne sur un grand ensemble de chemins de même variance, mais elles diffèrent sur chaque chemin pris isolément.}
\end{solution}

\begin{solution}{iq:iv2:crypto-for-traders-and-researchers:25}
Elle multiplie le capital par environ 20.4. Aux bornes haute et basse de l'intervalle, elle perd environ 2.32~\% et 2.83~\% par rapport à la détention, et la pleine largeur perd 0.11~\% et 0.14~\%. Au-delà de l'intervalle, elle ne détient qu'un seul actif et ne perçoit aucun frais.

\textbf{Capital.} Pour une liquidité $L$, la pleine largeur requiert une valeur de $2L\sqrt{P}$, et l'intervalle requiert $L(2\sqrt{P} - \sqrt{a} - P/\sqrt{b})$. Leur rapport, pour $a = 1\,800$, $b = 2\,200$, $P = 2\,000$, vaut 20.4.

\textbf{Perte.} Pour la position sur l'intervalle, on compare $V(P_1) = L(2\sqrt{P_1} - \sqrt{a} - P_1/\sqrt{b})$ à ce que vaudraient ses avoirs initiaux. À $P_1 = 2\,200$, le rapport vaut $-2.32~\%$. Pour la pleine largeur, il vaut $2\sqrt{1.1}/2.1 - 1 = -0.11~\%$. La perte par dollar augmente donc du même facteur, d'environ 20, de même que le taux de loss-versus-rebalancing et les revenus de frais.

\textbf{Au-delà de l'intervalle.} Au-dessus de 2\,200, la position ne détient que des dollars et manque toute hausse supplémentaire. En dessous de 1\,800, elle ne détient que de l'ETH et subit toute la baisse. Les frais s'arrêtent dans les deux cas.

La concentration amplifie donc tout avec l'effet de levier. C'est comparable à un strangle vendu sur ces strikes, et ce n'est pas une hausse gratuite du rendement.

\iqwrong{Affirmer que «~une efficacité du capital de 20x signifie un rendement 20 fois supérieur.~» Il s'agit de 20 fois les frais et de 20 fois le coût de gamma, sur une position qui peut sortir de son intervalle et cesser de rapporter.}
\end{solution}

\begin{solution}{iq:iv2:crypto-for-traders-and-researchers:26}
L'attaquant peut acheter environ 0.47 ETH et gagne environ 216~\$. Le surcoût de la victime est d'environ 223~\$. L'écart revient aux frais du pool. Regrouper les transactions à un prix unique, ou passer par un flux d'ordres privé, la rendrait inoffensive.

Le cours affiché à la victime est 22\,289, donc la limite est 22\,512. L'attaquant achète d'abord $a$, et le coût de la victime augmente avec $a$. On résout pour que le coût atteigne 22\,512~: $a = 0.470$.

L'attaquant paie 947.60, la victime paie 22\,511.98, et l'attaquant revend pour 1\,163.48. Son profit est donc de 215.88. Le surcoût pour la victime est $22\,511.98 - 22\,289.09 = 222.89$. Et l'écart d'environ 7 revient aux frais et aux LP.

Le profit croît avec la tolérance. À une limite de 0.05~\%, il n'est que d'environ 11~\$, et cela peut ne pas couvrir le coût de la transaction. De plus, le coût de priorité des transactions grignote sa marge, et des concurrents l'enchérissent auprès des validateurs.

Versions inoffensives~: enchères par lots avec un prix de compensation unique, routes de transaction privées, ou limite de glissement plus serrée, bien que celle-ci risque de faire échouer des transactions.

\iqwrong{Croire que la limite de glissement est une garantie de prix. Toute tolérance est un cadeau~: elle indique à l'attaquant combien il peut prendre.}
\end{solution}

\begin{solution}{iq:iv2:crypto-for-traders-and-researchers:27}
Cela coûte environ 0.0913 bitcoin, soit environ 5\,478~\$. Le delta en pièce est d'environ 0.454, contre 0.546 pour l'option en dollars. Par parité en pièce, $C - P = (S - K)/S = 0$ à la monnaie.

Le gain en pièce multiplié par $S$ vaut $(S - K)^+$, c'est donc la même option en dollars, et $C_{\$} = S\Phi(d_1) - K\Phi(d_2)$. Avec $d_1 = -d_2 = 0.115$~: $C_{\$} = 5\,478$. D'où $C_{\text{coin}} = 0.0913$. Une règle empirique~: $0.4 \times 0.8 \times \sqrt{30/365} = 0.0917$.

Le delta en dollars vaut $\Phi(d_1) = 0.546$. Le cours en pièce est $C_{\$}/S$. La sensibilité du cours en pièce, multipliée par $S$, vaut donc $\Phi(d_1) - C_{\$}/S = 0.454$. L'écart correspond à la prime~: on l'a payée en pièce, et cette pièce a elle-même un delta. Pour couvrir en dollars, il faut aussi compter la prime en pièce ou la marge.

\iqwrong{Couvrir à 0.546 sans corriger la prime lue en pièces. On serait alors surcouvert de l'exposition en pièce de la prime, si celle-ci est détenue en pièce.}
\end{solution}

\begin{solution}{iq:iv2:crypto-for-traders-and-researchers:28}
Le rendement net sur capital est d'environ 21~\%. Avec le risque de plateforme, il tombe à environ 18.5~\%, et le ratio de Sharpe vaut environ 1.7, et non 10. Le test porte en outre sur un an, et ce n'est pas la bonne année.

\textbf{Coûts.} $20 \text{ pb} \times 2 \text{ jambes} \times 12 = 4.8~\%$ du notionnel. Donc $30 - 4.8 = 25.2~\%$ sur le notionnel. Avec une marge de 20~\%, le capital vaut 1.2 fois le notionnel, donc $25.2/1.2 = 21~\%$.

\textbf{Plateforme.} Une probabilité de 5~\% de perdre la moitié coûte 2.5~\%, soit une moyenne de 18.5~\%. L'écart type de cet événement vaut $0.5\sqrt{0.05 \times 0.95} = 10.9~\%$. Avec 2~\% de bruit supplémentaire, le total fait 11.1~\%. Donc $18.5/11.1 = 1.67$.

\textbf{Le reste.} Un test sur un an sans mois perdant ne contient aucun événement du type de ceux qui tuent cette stratégie. Un Sharpe élevé qui cache une queue épaisse, c'est justement le problème.

\begin{dated}{2026-11}{Comment le funding évolue}\label{dat:iv2:crypto-for-traders-and-researchers:hist}
Sur une grande plateforme, le funding annualisé des perpetuals bitcoin s'est établi en moyenne à environ 17~\% en 2020, 31~\% en 2021, 4~\% en 2022, 8~\% en 2023, 12~\% en 2024, 5~\% en 2025 et 3~\% en 2026 jusqu'en septembre. En 2022, environ un cinquième des relevés étaient négatifs. Le test à 30~\% ne porte donc que sur une année sur six, et une année isolée dit peu de chose.
\end{dated}

Le reste de la liste~: un échantillon pris dans une année très active~; des jambes de couverture supposées exécutées intégralement~; des frais au niveau du preneur de liquidité sous-estimés~; des appels de marge non modélisés~; des plateformes ayant fait défaut par la suite exclues~; la capacité (le funding baisse à mesure que le capital arrive)~; et un ratio de Sharpe annoncé, mesuré sur des valorisations quotidiennes où la queue n'apparaît pas.

\iqwrong{Tester le chiffre et accepter la stratégie. Il faut se demander ce qui devrait se produire pour que cela perde la moitié du capital, et si cela figure dans les données. Ici, ce n'est pas le cas.}
\end{solution}

\begin{solution}{iq:iv2:crypto-for-traders-and-researchers:29}
À 2x, 5x et 10x, le rendement brut vaut 6.7~\%, 8.3~\% et 9.1~\%. Les probabilités de liquidation mensuelle sont de 7.7~\%, 42.7~\% et 67.8~\%. Les rendements nets valent environ 6.0~\%, 4.1~\% et 1.7~\%. Il faut un faible levier, peut-être environ 2x, à moins de pouvoir déplacer la marge rapidement.

\textbf{Brut.} Le capital achète $N$ de spot et $N/L$ de marge, donc le notionnel vaut $N = C/(1 + 1/L)$. À $L = 2$, il vaut $0.667C$, soit $10~\% \times 0.667 = 6.67~\%$. À 5, $8.33~\%$. À 10, $9.09~\%$.

\textbf{Liquidation.} La distance à la barrière vaut $1/L$ de la position vendeuse. On a donc $\ln 1.5$, $\ln 1.2$ et $\ln 1.1$, comme calculé avec $2\Phi(-b/(\sigma\sqrt{T}))$ et $T = 30/365$.

\textbf{Net.} Chaque liquidation coûte 1~\% du notionnel. Douze probabilités mensuelles de $p$ donnent $12p \times 1~\% \times N/C$~:
\begin{itemize}
\item 2x~: $0.925~\% \times 0.667 = 0.62~\%$, soit 6.05~\%.
\item 5x~: $5.12~\% \times 0.833 = 4.27~\%$, soit 4.07~\%.
\item 10x~: $8.13~\% \times 0.909 = 7.39~\%$, soit 1.70~\%.
\end{itemize}
Chercher donc du levier pour gagner 2.4~\% de rendement brut en plus coûte 6.8~\% de pertes supplémentaires en liquidation. Un coût de 1~\% est aussi un peu léger pour un marché violent, ce qui pousse encore davantage vers un faible levier.

\textbf{Autres points, avant de le lancer.}
\begin{itemize}
\item \textbf{Régime de financement.} 10~\% est une moyenne. Prévoyez une règle de sortie lorsque le financement reste inférieur à vos coûts. Les backtests doivent inclure les années à faible financement.
\item \textbf{Limites par plateforme.} Plafonner chaque plateforme. Savoir ce qui arrive à l'autre jambe si l'une tombe en panne, et prévoir la sortie à l'avance.
\item \textbf{Mécanique des marges.} La marge croisée sur une seule plateforme est nettement plus sûre, si l'on peut tolérer la concentration. Sinon, il faut un plan de transfert plus rapide que le mouvement, et du capital de réserve.
\item \textbf{Auto-deleveraging.} Votre vente à découvert rentable peut être réduite.
\item \textbf{Stablecoin.} Si vos dollars sont des stablecoins, vous portez aussi le risque de leur arrimage.
\item \textbf{Reporting.} Suivre séparément le carry, les frais, le glissement de la couverture et les pertes de queue, afin de voir quand l'avantage a disparu.
\end{itemize}

\iqwrong{Choisir le rendement le plus élevé. Un basis trade semble sans risque jusqu'à ce que le levier transforme un mouvement de prix en perte réalisée~; un levier élevé sur une transaction d'apparence sans risque est donc la façon dont ce type de desk disparaît.}
\end{solution}

\begin{solution}{iq:iv2:crypto-for-traders-and-researchers:30}
Sans frais, vous perdez environ 7.8~\%, soit 780\,000~\$. Les frais d'arbitrage représentent environ 0.7~\%, 2.9~\% et 5.2~\%. Pour atteindre le seuil de rentabilité, il faut un volume annuel de flux non informé d'environ 1.4~milliard de dollars à 5 points de base, 163~millions de dollars à 30, et 26~millions de dollars à 100. Je ne le prendrais que si je peux démontrer ce flux, car des frais élevés le font fuir, et l'arbitrage seul ne rapporte jamais.

\textbf{Simulation.} Prix horaires à 80~\% de volatilité, un arbitragiste déplaçant le pool jusqu'au bord de la bande de frais à chaque heure, et une couverture du delta du pool. Le net est ce qu'obtient un LP couvert en delta sans autre flux.

\begin{center}
\begin{tabular}{rrrrr}
\toprule
Frais (pb) & Perte & Frais d'arbitrage & Net & Volume requis (\$m/an) \\
\midrule
0 & 7.8~\% & 0 & $-7.8~\%$ & --- \\
5 & 7.8~\% & 0.7~\% & $-7.1~\%$ & 1\,430 \\
30 & 7.8~\% & 2.9~\% & $-4.9~\%$ & 163 \\
100 & 7.8~\% & 5.2~\% & $-2.6~\%$ & 26 \\
\bottomrule
\end{tabular}
\end{center}

La perte sans frais correspond à la théorie~: $1 - e^{-\sigma^2/8} = 7.7~\%$.

Les frais d'arbitrage seuls sont toujours inférieurs à la perte, et il doit en être ainsi~: l'arbitragiste n'échange que s'il gagne, et le résultat d'un LP couvert est l'opposé de ce gain. Reste donc ce que doit payer le flux non informé. À 30 points de base, $\$10 \text{ million} \times 4.9~\% = \$489\,000$, et en divisant par 0.3~\% on obtient 163~millions de dollars, soit environ 447\,000~\$ par jour.

\textbf{Dois-je le prendre~?} Je regarderais le flux non informé réel en part du pool, par niveau de frais. Des frais plus élevés demandent beaucoup moins de flux, mais ils font aussi fuir le flux vers des pools moins chers. La bonne question est le volume que j'obtiendrais réellement à chaque niveau de frais.

\textbf{Volatilité plus faible.} À 40~\% de volatilité, la perte tombe au quart, soit environ 2~\%, puisqu'elle dépend de $\sigma^2$. C'est donc une position de vente de volatilité, et je la comparerais à la vente directe d'options.

\textbf{Concentré.} Avec une plage de 20x, les frais, le gamma et la perte sont multipliés par 20 dans la plage, et la position peut sortir de la plage. Le flux de seuil de rentabilité par unité de liquidité reste à peu près le même, donc le flux nécessaire par dollar de capital est environ 20 fois plus grand, et ce flux est partagé avec d'autres LP. Et le risque de rester hors de la plage après un saut est nouveau.

\textbf{Pratique.} La couverture coûte le funding perpétuel sur une vente à découvert, elle a des frais, et elle ne peut pas être continue. Chacun de ces éléments s'ajoute à la perte. Et le risque de smart contract, ainsi que la jambe stablecoin, ne figurent pas dans le tableau.

\iqwrong{Comparer les frais perçus avec la perte impermanente sur une période où le prix a fini là où il avait commencé, et en conclure que les LP disposent d'un rendement gratuit. Un LP couvert paie la variance chaque jour, quel que soit le prix final.}
\end{solution}
